\textit{[ This chapter includes parts of a paper under review as - Alexander Politowicz, Sahisnu Mazumder, and Bing Liu. “Achieving State-of-the-Art Out-of-Distribution Detection Performance with only LLM Prompting”. Submitted to International Joint Conference on Artificial Intelligence (IJCAI 2026), 2026.]}

\section{Importance of Prompt-Only OOD Detection}
\label{sec:propwork-atpo}

Continuing the theme in this dissertation of improving ML for open-world application via OOD detection, advancing the accessibility of OOD detection methods is critical for enabling the usage of robust and effective ML in real-world systems. With the prevalence of large language models (LLMs) in modern society and technology, this factor of accessibility is now a paramount issue. Development of OOD detection on (and for) LLMs helps to not only reduce the issues inherent in LLMs, e.g. hallucination or incorrect reasoning \cite{sriramanan2024llm}, but also build trust in (and reliability of) LLMs deployed in production settings, products, and services \cite{liu2023good,ouyang2025textual}. However, the current landscape of OOD detection in LLMs is impenetrable for non-expert users. All previous OOD detection methods designed for usage in NLP and LLMs, except one, either rely on access to model internal components and information or require training of the model. To a non-expert user attempting to deploy such methods on their system, these requirements are intractable and OOD detection is likely to be discarded. Instead, the goal is to design OOD detection methods around an accessible interface which is not only already intuitive for users, that is via \textit{prompts}, but is also able to induce effective OOD detection in the LLM being used. This allows for the usage of OOD detection at a negligible cognitive and technological cost to the user.

%Solving the problem of out-of-distribution (OOD) detection with large language models (LLMs) using only natural language prompts is challenging. Without access to the model features and logits traditionally used in OOD detection, prompt-only methods must leverage the knowledge embedded within LLMs or provided with in-context examples. However, the previous in-context prompt-only OOD detection method used few-shot ($C$+1)-classification which is difficult for LLMs to solve when test samples lie near the decision boundary (denoted as ``near-OOD'' cases). In this work, we instead propose prompting the LLM to use \textit{lexical and semantic alignment} and \textit{OOD score-based thresholding} to perform OOD detection. This directly leverages LLM knowledge and provided in-context examples to generate a rich OOD score space for OOD detection \textit{without relying on fine-tuning or model features and logits}. We demonstrate that our method, \textbf{A}lignment-based \textbf{T}hresholding for \textbf{P}rompt-only \textbf{O}OD detection (\sys), outperforms the previous prompt-only OOD detection method and can even surpass strong traditional OOD detection approaches across near-OOD datasets, LLM families, model sizes, and few-shot example set sizes.

%Machine learning (ML) models typically operate under the \textit{closed-world assumption}~\citep{fei2016breaking,shu2017doc,joseph2021towards}, expecting test data to follow the same distribution as training data (\textit{in-distribution (ID)} data). However, real-world deployment often encounters \textit{out-of-distribution (OOD)} samples that differ semantically from training data, such as new classes with different label spaces. OOD detection identifies these semantic distribution shifts (e.g., due to \textit{new class occurrences}) and prevents models from making unreliable predictions. This capability is crucial for deploying many real-world applications such as dialogue intent detection \citep{lang2024out} and text classification \citep{baran2023classical}, enhancing model reliability and trustworthiness in dynamic environments where novel inputs are common \citep{fei2016breaking,zhang2023openood,kaymak2025learning}. 

%A large body of work \citep{yang2021generalized,lu2025out} exists on OOD detection, both in vision and text domains. Our work focuses on text \citep{lang2023survey,xu2025large}. Existing OOD detection approaches all require fine-tuning Pre-trained Language Models (PLMs) on ID data or using the features and logits of PLMs with additional sophisticated techniques \citep{salimbeni2024beyond,lee2024ced}. Developing OOD detectors thus involved significant manual effort (e.g., labeling a large amount of ID training data, ML software programming, and resource-intensive model training) as well as technical field expertise. More recently, with the advent of powerful LLMs capable of solving a wide-range of complex natural language tasks through reasoning, one work \citep{wang2024beyond} has introduced the problem of \textbf{prompt-only OOD detection}. The goal is to leverage an existing LLM to build an OOD detector \textit{through only prompting and In-Context Learning (ICL)} without requiring any expensive model training or fine-tuning, large scale data collection, or sophisticated ML programming. This significantly reduces model development costs and manual effort. Specifically, \citet{wang2024beyond} studied OOD detection as an ($C$+1)-classification problem (involving $C$ ID classes and one ``\textit{unknown}" class representing OOD) \citep{shu2017doc} and proposed a few-shot ICL method that assumes no access to model features or logits. The method prompts an LLM with ID class examples and the input text (test sample) and parses the generated output for either an ID class or ``\textit{unknown}'' if no ID class was relevant.
To be specific, there exist many LLMs that are pre-trained on some corpus of data (also referred to as Pre-trained Language Models, or PLMs). OOD detection in the text domain has recently focused on leveraging these models by fine-tuning them on ID data or using the features and logits from the LLM with some additional technique for detection at test-time (see Chapter \ref{diss_relwork}). These approaches require significant manual effort. Fine-tuning requires users to label large amounts of ID training data, have the expertise to program the ML software and handle the models involved, and obtain access to expensive computational resources for the training procedures. Feature- and logits-based approaches, of which have been discussed in this thesis extensively, assume the user understands the fundamental components of LLMs used and that they can program routines which both access and effectively leverage the requisite features and logits. Neither of these approaches is reasonable for any non-expert user and thus establishes an impossible barrier between real-world application of ML models and OOD detection.

With the advent of reasoning in LLMs, and its demonstrated ability to solve complex tasks via natural language and prompting \cite{kumar2025llm,zhang2022automatic,dong2024survey}, one previous work in the literature has considered the approach of \textit{prompt-only OOD detection} \cite{wang2024beyond}. This approach uniquely limits the OOD detection method to only a prompt and minimal post-processing of the LLM's response. The goal is to utilize an existing LLM to build an OOD detector using only prompting, without any model training, fine-tuning, access to features and/or logits, data collection, or ML programming. Eliminating these factors bridges the gap between LLMs and OOD detection, a major step forward for the problem of real-world usage of ML.

However, the method is significantly limited. The authors propose to approach the OOD detection problem as a ($C$+1)-classification problem, where the targets considered by the LLM are the $C$ ID classes and one ``unknown'' class (representing the entirety of the OOD data). They use few-shot In-Context Learning \cite{dong2024survey} to investigate the performance of various LLMs across various datasets using only prompts containing ID class examples and the input test utterance. Output was parsed for either the ID class or the exact text of ``unknown'' given that no ID class was considered relevant by the model.

%We argue that this simple classification-based prompting for OOD detection is sub-optimal, which explains its poor performance (see Sec \ref{sec:results}). While ($C$+1)-class decision boundaries work well when ID and OOD samples are well-separated in the embedding space, they fail in near-OOD scenarios \citep{zhang2023openood}, where ID and OOD data lie close together or are inseparable (see Section \ref{sec:exp}). In such complex cases, finding appropriate decision boundaries to distinguish ID from OOD data is crucial for OOD detection performance. Direct ($C$+1)-classification LLM prompting as used in \cite{wang2024beyond} often fails because LLMs cannot effectively discriminate boundaries using only a few examples, i.e., without considering the full data distribution. Since existing feature-based OOD methods leverage the complete ID data distribution during inference, they continue to outperform prompt-only approaches in OOD detection tasks.
The previous proposed approach to prompt-only OOD detection, while novel and establishing the task, is sub-optimal. In particular, ($C$+1)-classification approaches struggle when ID and OOD samples lie close together (or are inseparable) relative to decision boundaries in the embedding space. Such near-OOD samples cannot be easily distinguished using decision boundaries as ID or OOD. However, it is exactly such cases where OOD detection is most critical. LLMs cannot effectively discriminate these boundaries separating ID and near-OOD data primarily because they must rely on reasoning and generation from only a few ICL examples. Traditional OOD detection methods have access to the entire ID data distribution and can leverage the information within the distribution during OOD detection to more completely compare potential OOD samples to seen ID data behaviors. It is for this reason that these traditional OOD detection methods continue to outperform prompt-only OOD detection.

%In this paper, we devise a different prompt-only approach for OOD detection and propose a novel solution that overcomes the limitations of the existing prompt-only method. We prompt LLM to 1) perform alignment between the input text and ICL context (including ID class categories and few-shot examples for each class), and then use the alignment information to 2) generate an OOD score that estimates the degree of distribution shift between the input and the ID class categories. The estimated scores for the ID data distribution are used to calculate an ID-informed threshold for OOD detection. Our method, \textbf{Alignment-based Thresholding for Prompt-only OOD detection (\sys)}, not only outperforms the previous prompt-only OOD detection method \citep{wang2024beyond}, which, to our knowledge, is the only existing prompt-only OOD detection method, by a large margin but also improves over traditional OOD detection methods with access to model feature and logits. Our experiments across a range of near-OOD datasets, LLM families, model sizes, and few-shot example configurations demonstrates the contribution of \sys.
%\alex{In this paper, we instead propose an advancement to In-context Prompt-only OOD Detection (which we denote as \sys) and show that it can achieve the same or even better OOD detection performance using in-context learning (ICL) \textbf{with only \textit{verbal prompts} and \textit{no coding}}. This is of great importance because it means that anyone with no programming knowledge can perform the task.
Overcoming this issue requires a re-framing of the prompt-only OOD detection approach. Instead of limiting the LLM to ($C$+1)-classification, leveraging OOD scores, as traditional OOD detection methods do, enables a greater capacity for both reasoning and detection via reasoning and referencing of the ID data distribution. Specifically, the proposed method, \textbf{A}lignment-based \textbf{T}hresholding for \textbf{P}rompt-only \textbf{O}OD detection (\syst), uses 1) explicitly induced alignment between the input test utterance and the ICL context (composed of the ID classes and few-shot examples for each class), and 2) generation of an OOD score estimating the degree of shift between the input and the ID data distribution. Scores for the ID data can be easily collected and used (during post-processing, preserving accessibility) to calculate an ID data distribution-informed threshold which can subsequently be used for OOD detection. This approach builds a rich score space for informed OOD detection, supported by alignment-based reasoning, instead of relying on the classification of a single sample with no reasoning or space for consideration of uncertain samples.

%The goal is to generate low OOD scores $S$ for ID test data $\mathcal{D}_{ID}$ and high OOD scores $S$ for OOD test data $\mathcal{D}_{OOD}$.
% Our experiments across a range of text datasets, LLMs, and verbalized OOD methods verify that \sys significantly improves prompt-only OOD detection over existing prompting methods. {\color{blue}We further investigate prompt-only OOD detection methods across a range of prompt and ICL variations. Our findings show that prompt-only OOD detection performance is tied to the reasoning capabilities of the LLM and alignment between a test sample and the given ICL context (see Section TODO).}

%\textbf{Prompt-Only OOD Detection with LLMs.} 
%Recently, \citet{wang2024beyond} proposed a prompt-only ($C$+1)-classification method using LLMs, assuming no access to model features or logits. It addresses the OOD detection task tangentially but is the \textbf{only} previous work that considers prompt-only OOD detection.
%Specifically, it prompts an LLM with ID class examples and the input query before model output is parsed for either an ID class or ``unknown'' (OOD) if no ID class was relevant. 
%
%% The LLM output was parsed to identify if the correct class was generated for ID test samples or ``unknown'' was generated for OOD samples. 
%%To our knowledge, it is the sole study in this direction. However, our experiments show that our \sys outperforms this method. %We advance this direction by proposing to convert existing OOD detection methods into a prompt. % demonstrate that using more intelligent reasoning from the LLM, state-of-the-art prompt-only OOD detection can be achieved by 
%%It converts an existing OOD detection technique into a prompt.
%We evaluated their method both in the open-set classification setting and the OOD detection setting, and showed that in both cases the approach was suboptimal. Instead, we show that leveraging alignment with LLM reasoning and thesholding of OOD scores greatly improves performance and, interestingly, can even outperform feature- and logits-based OOD detection methods. We induce alignment through prompting of an LLM and find the utilization of both lexical alignment \citep{el2021xlent,dyer2013simple,srivastava2025measuring} and semantic alignment \citep{el2021xlent,thompson2020cultural} to be beneficial. Previous work on alignment manually defines alignment metrics and do not consider ICL. %, while we induce alignment in the generation of an LLM.
%More related work is in Appendix Section 3.
Evaluation of ATPO is conducted, as in all previous works, over a range of near-OOD datasets both against the only previous prompt-only OOD detection method \cite{wang2024beyond} and against traditional OOD detection methods with access to features and logits. Datasets span various application domains, including intent detection (Banking \cite{casanueva2020efficient}) and topic classification (20NewsGroups \cite{lang1995newsweeder}), and exclusively involve ID-OOD data splits from within the same dataset, establishing potent near-OOD settings with significant lexical, semantic, topical, and covariate overlap. Extensive ablation studies further explore the source of ATPO performance and identify the key components necessary for prompt-only OOD detection, significantly advancing it into the frame of OOD detection and establishing it as a state-of-the-art, accessible OOD detection method for modern techonological solutions.
